\section*{Розділ IV. Голова ТПК}
\addcontentsline{toc}{section}{Розділ IV. Голова ТПК}

\subsection*{4.1. Призначення та припинення повноважень}
\addcontentsline{toc}{subsection}{4.1. Призначення та припинення повноважень}
    4.1.1. Голова ТПК призначається рішенням КСУ.

    4.1.2. Головою ТПК може бути студент Університету денної форми навчання.

    4.1.3. Повноваження Голови ТПК припиняються у разі:

        \begin{enumerate}[label=\alph*)]
            \item Відкликання рішенням КСУ;
            \item Подання особистої заяви про складення повноважень на ім'я Голови КСУ;
            \item Втрати статусу студента Університету денної форми навчання.
        \end{enumerate}

    4.1.4. У разі дострокового припинення повноважень Голови ТПК, КСУ призначає нового Голову ТПК на найближчому засіданні. До призначення нового Голови ТПК його обов'язки виконує Голова СР ТПК КАІ.

\subsection*{4.2. Повноваження Голови ТПК}
\addcontentsline{toc}{subsection}{4.2. Повноваження Голови ТПК}
    4.2.1. Голова ТПК:

        \begin{enumerate}[label=\alph*)]
            \item Здійснює загальне керівництво діяльністю ТПК;
            \item Представляє ТПК у відносинах з адміністрацією Університету, іншими органами та організаціями;
            \item Підписує документи ТПК;
            \item Звітує перед КСУ про діяльність ТПК;
            \item Видає розпорядження з питань діяльності ТПК;
            \item Координує діяльність Голови СР ТПК КАІ та Голови СР ТПК СМ;
            \item Здійснює інші повноваження, передбачені цим Положенням та рішеннями КСУ.
        \end{enumerate}

\subsection*{4.3. Колегіальне прийняття рішень}
\addcontentsline{toc}{subsection}{4.3. Колегіальне прийняття рішень}
    4.3.1. Наступні рішення приймаються колегіально Головою ТПК, Головою СР ТПК КАІ та Головою СР ТПК СМ більшістю голосів:

        \begin{enumerate}[label=\alph*)]
            \item Призначення та зняття голів СР ТПК факультетів/інститутів та СР ТПК гуртожитків.
        \end{enumerate}

    4.3.2. Ініціювати колегіальний розгляд питання може будь-хто з трьох посадових осіб, зазначених у п. 4.3.1.

    4.3.3. Рішення за результатами колегіального розгляду оформлюються розпорядженням Голови ТПК.
